% =============================================================================
\subsection{Tier 0: Self-supervised pretraining}
\label{sec:tier0}
% =============================================================================

\begin{figure}[t]
  \centering
  % =============================================================================
% Tier 0 SSL pipeline — detailed figure for the paper.
%
% A standalone, self-contained TikZ figure showing the flow of data and
% gradients through the self-supervised pretraining stage. Layout is a
% left-to-right pipeline with four columns:
%
%   [raw -> mask -> masked]  ->  [ASTA encoder]  ->  [SSL head]  ->  [loss]
%                                       |                               ^
%                                       v                               |
%                              [pretrained encoder]  (reused downstream)
%
% Compile via \input{} from main.tex (which already loads tikz), or
% standalone with:
%   pdflatex tier0_pipeline.tex
%
% Style matches figures/architecture.tex so the figures feel
% consistent: same sffamily font, same rounded-corner blocks, same
% blue/orange/teal palette.
% =============================================================================
\begin{tikzpicture}[
  font=\sffamily\small,
  node distance=8mm and 6mm,
  every node/.style={transform shape},
  cell/.style={draw, minimum width=2.4mm, minimum height=2.4mm, inner sep=0pt,
                fill=blue!8},
  cellaoa/.style={draw, minimum width=2.4mm, minimum height=2.4mm, inner sep=0pt,
                   fill=orange!15},
  cellmask/.style={draw, minimum width=2.4mm, minimum height=2.4mm, inner sep=0pt,
                   fill=blue!4},
  cellmaskaoa/.style={draw, minimum width=2.4mm, minimum height=2.4mm, inner sep=0pt,
                      fill=orange!4},
  cellblk/.style={draw, minimum width=2.4mm, minimum height=2.4mm, inner sep=0pt,
                  fill=black!90},
  block/.style={draw, rounded corners=2pt, minimum width=18mm,
                minimum height=9mm, align=center, fill=blue!10},
  acc/.style={draw, rounded corners=2pt, minimum width=18mm,
              minimum height=9mm, align=center, fill=orange!18},
  feat/.style={draw, rounded corners=2pt, minimum width=18mm,
               minimum height=9mm, align=center, fill=teal!18},
  in/.style={draw, rounded corners=2pt, minimum width=18mm,
             minimum height=9mm, align=center, fill=blue!22},
  out/.style={draw, rounded corners=2pt, minimum width=24mm,
              minimum height=10mm, align=center, fill=teal!22},
  loss/.style={draw, rounded corners=2pt, minimum width=22mm,
               minimum height=10mm, align=center, fill=red!14},
  arr/.style={->, >=stealth, thick},
  arrdash/.style={->, >=stealth, thick, dashed, color=blue!50!black},
  arrthin/.style={->, >=stealth, thin, color=blue!50!black},
  arrupd/.style={->, >=stealth, thick},
]

% =============================================================================
% COLUMN 1: data prep (raw -> mask -> masked)
% =============================================================================
\begin{scope}[local bounding box=col1, shift={(0,0)}]
  % --- raw window (ToA on top, AoA below) ---
  \node[anchor=south] at (10mm, 35mm) {\textbf{raw window} $\bm{X}$};
  \foreach \t in {0,...,7} {
    \foreach \m in {0,...,3} {
      \node[cell, minimum width=2.4mm, minimum height=2.4mm, inner sep=0pt]
        at (\t*2.4mm + 1.5mm, 32mm - \m*2.4mm) {};
      \node[cellaoa, minimum width=2.4mm, minimum height=2.4mm, inner sep=0pt]
        at (\t*2.4mm + 1.5mm, 21mm - \m*2.4mm) {};
    }
  }
  \node[anchor=west, font=\scriptsize, color=blue!60!black]
    at (22mm, 28mm) {ch 0: ToA (m)};
  \node[anchor=west, font=\scriptsize, color=orange!70!black]
    at (22mm, 17mm) {ch 1: AoA (°)};
  \node[anchor=west, font=\scriptsize, color=blue!60!black]
    at (22mm, 11mm) {$(T{=}32,\,M{=}4,\,F{=}2)$};

  % --- random mask ---
  \node[anchor=south] at (10mm, 8mm) {\textbf{random mask} $\bm{M}$};
  \foreach \t/\m/\b in {
    0/1/1, 0/3/1,
    1/2/1,
    2/2/1,
    3/0/1,
    4/1/1,
    5/2/1,
    6/2/1,
    7/0/1, 7/3/1} {
    \ifnum\b=1
      \node[cellblk] at (\t*2.4mm + 1.5mm, 5mm - \m*2.4mm) {};
    \else
      \node[cellmask] at (\t*2.4mm + 1.5mm, 5mm - \m*2.4mm) {};
    \fi
  }
  \node[anchor=west, font=\scriptsize, color=blue!60!black]
    at (22mm, 1mm) {$(T,M)$ bool, $p{=}0.25$};
  \node[anchor=west, font=\scriptsize, color=red!50!black]
    at (22mm, -2mm) {joint across $F$};

  % --- masked input ---
  \node[anchor=south] at (10mm, -8mm) {\textbf{masked input} $\tilde{\bm{X}}$};
  \foreach \t/\m/\b in {
    0/1/1, 0/3/1,
    1/2/1,
    2/2/1,
    3/0/1,
    4/1/1,
    5/2/1,
    6/2/1,
    7/0/1, 7/3/1} {
    \ifnum\b=1
      \node[cellblk] at (\t*2.4mm + 1.5mm, -11mm - \m*2.4mm) {};
      \node[cellblk] at (\t*2.4mm + 1.5mm, -22mm - \m*2.4mm) {};
    \else
      \node[cell] at (\t*2.4mm + 1.5mm, -11mm - \m*2.4mm) {};
      \node[cellaoa] at (\t*2.4mm + 1.5mm, -22mm - \m*2.4mm) {};
    \fi
  }

  % raw -> mask
  \draw[arr] (10mm, 19mm) -- (10mm, 11mm)
    node[midway, right, font=\scriptsize, color=blue!60!black] {sample};
  % mask -> masked
  \draw[arr] (10mm, -1mm) -- (10mm, -7mm)
    node[midway, right, font=\scriptsize, color=blue!60!black] {apply};
\end{scope}

% encode arrow: from right of masked input, right then up to TCN
\begin{scope}
  \coordinate (enc_start) at (22mm, -16mm);   % right of masked input, mid AoA row
  \coordinate (enc_mid)   at (35mm, -16mm);   % midway horizontally
  \coordinate (enc_top)   at (35mm, 32mm);     % up to TCN level
  \draw[arr] (enc_start) -- (enc_mid);
  \draw[arr] (enc_mid) -- (enc_top);
  \node[font=\scriptsize, color=blue!60!black, fontstyle=\itshape]
    at (35mm, -19mm) {encode};
\end{scope}

% =============================================================================
% COLUMN 2: ASTA encoder
% =============================================================================
\begin{scope}[local bounding box=col2, shift={(50mm, 0)}]
  \node[anchor=south, font=\bfseries] at (0, 36mm) {ASTA encoder $f_\theta$};
  \node[anchor=north, font=\scriptsize, color=blue!60!black, fontstyle=\itshape]
    at (0, 36mm) {$(T,M,F) \to (T,M,d{=}64)$};

  \node[block, fill=blue!10, minimum width=20mm]
    (tcn) at (0, 32mm) {Causal TCN\\(per-channel, $k{=}3$)};
  \node[block, fill=teal!18, minimum width=20mm]
    (gcn) at (0, 25mm) {Dynamic GCN\\$A^{(c)}$ per-channel};
  \node[block, fill=orange!18, minimum width=20mm]
    (mhsa) at (0, 18mm) {Multi-Head\\Self-Attn\\$h{=}4$, $d{=}64$};

  \draw[arr] (tcn) -- (gcn) node[midway, right, font=\scriptsize] {$\bm{H}$};
  \draw[arr] (gcn) -- (mhsa);

  % residual bracket on the right
  \draw[arrthin] (10mm, 32mm) -- (10mm, 18mm);
  \node[anchor=west, font=\scriptsize, color=blue!50!black, fontstyle=\itshape]
    at (11mm, 25mm) {+ residual};
\end{scope}

% encoder -> head: short diagonal up-right
\begin{scope}
  \coordinate (enc_out) at (60mm, 18mm);
  \coordinate (head_in) at (82mm, 28mm);
  \draw[arr] (enc_out) -- (head_in)
    node[midway, above, font=\scriptsize, color=teal!60!black, fontstyle=\itshape]
    {$\bm{H}$};
\end{scope}

% =============================================================================
% COLUMN 3: SSL head
% =============================================================================
\begin{scope}[local bounding box=col3, shift={(82mm, 0)}]
  \node[anchor=south, font=\bfseries] at (0, 36mm) {SSL head $g$};
  \node[anchor=north, font=\scriptsize, color=blue!60!black, fontstyle=\itshape]
    at (0, 36mm) {2-layer MLP, $d \to d \to F$};

  \node[block, minimum width=20mm] (lin1) at (0, 28mm) {Linear$(d{\to}d)$\\$+$ GELU};
  \node[block, minimum width=20mm] (lin2) at (0, 21mm) {Linear$(d{\to}F)$};

  \draw[arr] (lin1) -- (lin2);

  \node[anchor=west, font=\scriptsize, color=orange!70!black, fontstyle=\itshape]
    at (12mm, 21mm) {$\hat{\bm{X}}$ (ToA+AoA)};
\end{scope}

% head -> loss: short diagonal up-right
\begin{scope}
  \coordinate (head_out) at (92mm, 21mm);
  \coordinate (loss_in)  at (113mm, 30mm);
  \draw[arr] (head_out) -- (loss_in)
    node[midway, above, font=\scriptsize, color=orange!70!black, fontstyle=\itshape]
    {$\hat{\bm{X}}$};
\end{scope}

% =============================================================================
% COLUMN 4: loss
% =============================================================================
\begin{scope}[local bounding box=col4, shift={(113mm, 0)}]
  \node[anchor=south, font=\bfseries] at (0, 36mm) {loss};

  \node[loss, minimum width=28mm, minimum height=11mm]
    (loss_eq) at (0, 30mm) {$\mathcal{L}{=}\frac{\sum \bm{M}\!\odot\!(\hat{\bm{X}}{-}\bm{X})^2}{\sum \bm{M}}$};

  \node[anchor=west, font=\scriptsize, color=red!50!black, fontstyle=\itshape, align=left]
    at (15mm, 28mm) {per-channel\\weights\\$w_{\text{ToA}}{=}3$};

  \node[loss, minimum width=22mm, minimum height=8mm]
    (adam) at (0, 19mm) {Adam, $\eta{=}10^{-3}$\\30--100 epochs};

  \draw[arr] (loss_eq) -- (adam);
\end{scope}

% =============================================================================
% TRUE-X dashed path: from raw window, over the top, down to loss
% =============================================================================
\begin{scope}
  \coordinate (true_start) at (10mm, 32mm);  % top of raw window
  \draw[arrdash] (true_start) -- (10mm, 45mm);
  \draw[arrdash] (10mm, 45mm) -- (113mm, 45mm);
  \draw[arrdash] (113mm, 45mm) -- (113mm, 35.5mm);
  \node[font=\scriptsize, color=blue!50!black, fontstyle=\itshape]
    at (60mm, 47mm) {true $\bm{X}$ (from data)};
\end{scope}

% =============================================================================
% BELOW column 2: pretrained encoder output box
% =============================================================================
\begin{scope}
  \node[out] (enc_out_box) at (50mm, 5mm)
    {\textbf{pretrained encoder} $f_{\theta^*}$\\(reused by Tier 1 / 3)};

  % parameters arrow from MHSA down to encoder output
  \draw[arr, color=teal!60!black] (50mm, 14mm) -- (50mm, 10mm)
    node[midway, right, font=\scriptsize, color=teal!60!black, fontstyle=\itshape]
    {params};

  % update arrow from Adam down/left to encoder output
  \draw[arrupd] (105mm, 15mm) .. controls (95mm, 8mm) .. (65mm, 6mm)
    node[pos=0.4, below, font=\scriptsize, fontstyle=\itshape] {update};
\end{scope}

\end{tikzpicture}

  \caption{Tier 0 self-supervised pretraining. A raw UWB window
  $\bm{X} \in \R^{T \times M \times F}$ is corrupted by a Bernoulli
  mask $\bm{M}$ (joint across channels) to produce the masked
  input $\tilde{\bm{X}}$. The ASTA encoder $f_\theta$ produces
  per-cell features $\bm{H}$, and a 2-layer SSL head $g$ predicts
  $\hat{\bm{X}}$ at every cell. The loss is the masked-MSE between
  $\hat{\bm{X}}$ and the original $\bm{X}$ on the masked cells only;
  for $F=2$ we weight the per-channel terms ($w_{\mathrm{ToA}}{=}3$,
  $w_{\mathrm{AoA}}{=}1$) so the geometrically meaningful ToA channel
  dominates the gradient. The pretrained encoder $f_{\theta^*}$ is
  the artifact passed to Tier 1 (FOMAML) and Tier 3 (downstream heads).}
  \label{fig:tier0}
\end{figure}

We pretrain the ASTA encoder with a self-supervised objective
that requires no labels. The motivation is twofold: (i) unlabeled
UWB streams are far cheaper to collect than labeled traces, and
(ii) decoupling the learning of \emph{measurement structure}
(channel-agnostic) from \emph{deployment structure}
(environment-specific, deferred to Tier~1) avoids the
under-supervision of the former during meta-learning.

\paragraph{Masked prediction.}
Given a batch $\{\bm{X}^{(b)}\}_{b=1}^{B}$, we sample a binary mask
$\bm{M} \in \{0,1\}^{T \times M}$ with entries drawn i.i.d.\ from
$\mathrm{Bernoulli}(p)$. The masked input $\tilde{\bm{X}}$
zeroes out the masked cells. The encoder produces
$\bm{H} = f_\theta(\tilde{\bm{X}})$, and a 2-layer MLP
$g_\eta : \R^d \to \R^F$ predicts
$\hat{\bm{X}} = g_\eta(\bm{H})$ at every (anchor, time) cell.
The self-supervised loss is the MSE on masked cells only:
\begin{equation}
  \mathcal{L}_{\mathrm{ssl}}(\theta, \eta)
  = \frac{1}{\sum_{b,t,m} \bm{M}^{(b)}_{t,m}}
    \sum_{b,t,m,f} \bm{M}^{(b)}_{t,m}
    \left( \hat{\bm{X}}^{(b)}_{t,m,f} - \bm{X}^{(b)}_{t,m,f} \right)^2.
  \label{eq:ssl}
\end{equation}
We use mask ratio $p = 0.25$. Augmentations (time-warp, anchor
dropout, additive Gaussian noise) are applied to the input before
masking. We considered a contrastive (NT-Xent) auxiliary loss
\cite{simclr2020} but found it interfered with masked prediction
during early training; we omit it.

\paragraph{Refill quality.}
After pretraining, we measure how well the encoder+head refills
masked cells on a held-out \emph{empty\_open} validation split
($F{=}2$, 25\% mask, $n{=}7{,}670$ masked cells per channel).
Table~\ref{tab:tier0_3to1_peranchor} reports per-anchor RMSE and
MAE for both channels under the selected \emph{3:1} (ToA:AoA)
channel weighting. The model achieves $0.151$\,m RMSE on ToA
($0.102$\,m MAE) and $5.21^\circ$ RMSE on AoA azimuth
($3.58^\circ$ MAE), with per-anchor spread of $\pm 1.0$\,cm on
ToA and $\pm 0.78^\circ$ on AoA — the refill behaves consistently
across the array, with no single anchor acting as a
disproportionately easy or hard case. Both channels operate
well above the visible-cell noise floor
($\sigma_{\text{ToA}}{=}0.050$\,m), confirming that the encoder
is learning nontrivial interpolation rather than a denoising
identity. The RMSE/MAE ratio of $\sim 1.5$ on both channels
indicates a moderately heavy-tailed error distribution
($\sqrt{\pi/2}{\approx}1.25$ for a clean Gaussian): the
occasional large miss is consistent with a small fraction of
geometrically degenerate windows where the visible-cell context
is too sparse to pin down the masked anchor.
\looseness=-1

% Tier-0 SSL — F=2 (refill), 25% mask, 3:1 (ToA:AoA) channel weighting
% Evaluated on Sionna `empty_open` validation split (240 windows,
% 7,670 masked cells). Reproduced from
% `code/checkpoints/tier0_empty_open_f2_25pct_3to1_100ep.pt` using
% `code/scripts/viz_tier0_metrics_f2.py` after fixing a double-masking
% bug in the eval pipeline (the model now masks internally; the
% outer eval-time `random_mask` was overwriting the skip-connection
% passthrough at visible cells).

\begin{table}[t]
\centering
\caption{Per-anchor Tier-0 SSL reconstruction error for the selected
\emph{3:1} (ToA:AoA) channel-weighted masked-cell loss, on the
Sionna \texttt{empty\_open} validation split ($F{=}2$, 25\% mask,
$n_{\text{masked}}{=}7{,}670$ cells per channel).
Per-anchor spread is tight on both channels: $\pm 1.0$\,cm on ToA
and $\pm 0.79^\circ$ on AoA. AoA error on A2 is again the lowest
because A2's azimuth is near-constant and trivially predicted;
the non-trivial-anchor mean $\overline{\text{RMSE}}_{\text{AoA,A1,A3,A4}}
= 5.54^\circ$ is reported as a fairer summary.}
\label{tab:tier0_3to1_peranchor}
\small
\begin{tabular}{ll rrrr r r}
\toprule
Channel & Metric & {A1} & {A2} & {A3} & {A4} & {$\bar{\cdot}$} & {Spread} \\
\midrule
ToA (m)         & RMSE & 0.151 & 0.163 & \textbf{0.143} & 0.148 & 0.151 & $\pm 0.010$ \\
ToA (m)         & MAE  & 0.103 & 0.105 & \textbf{0.100} & 0.101 & 0.102 & $\pm 0.003$ \\
\addlinespace
AoA azimuth     & RMSE & 5.158 & \textbf{4.225} & 5.655 & 5.791 & 5.207 & $\pm 0.78$ \\
AoA azimuth     & MAE  & 3.306 & \textbf{2.932} & 4.076 & 4.022 & 3.584 & $\pm 0.57$ \\
\bottomrule
\end{tabular}
\end{table}


\paragraph{Linear-probe diagnostic.}
We report a linear-probe localization error on a held-out
environment to verify that the pretrained encoder transfers
useful features. Section~\ref{sec:experiments} (in the
camera-ready) gives the numerical results.
